\section{Background / Literature Review（背景・先行研究）}

\subsection{観光地域における大型宿泊施設}

大型宿泊施設は観光地域における基幹インフラとして機能している。

特にリゾート型ホテルは、

\begin{itemize}
\item 宿泊
\item 飲食
\item 温泉
\item イベント
\item 集客
\end{itemize}

などを包括的に担い、観光地域全体の流動人口を支えている。

しかし、大型宿泊施設が持つ空間・集客力・文化発信力は、単なる商業機能だけではなく、地域文化を媒介する基盤としても活用可能である。

本研究では、大型宿泊施設を単なる宿泊機能としてではなく、「地域文化を接続する装置」として捉える。


\subsection{参加型観光研究}

近年、観光研究分野では、

\begin{itemize}
\item Participatory Tourism
\item Community-based Tourism
\item Experience Design
\end{itemize}

など、「参加」を重視する研究が進められている。

しかしその多くは、

\begin{itemize}
\item 地域体験消費
\item アクティビティ参加
\item 経済効果
\end{itemize}

を中心としており、

\begin{itemize}
\item 共に歌う
\item 同じ時間を過ごす
\item 感情を共有する
\end{itemize}

といった「共同体験構造」については十分に議論されていない。

本研究では、「共に時間を過ごすこと」そのものが観光価値となる構造に着目する。
